% ==========================================================
\section{Models@run.time}
\label{sec:bg_models_at_runtime}
% ==========================================================

\textit{Models@run.time} extend the principles of Model-Driven Engineering from the traditional development phase into the active execution environment of a software system.
By definition, they are causally connected self-representations of a software system that emphasize its structure, behavior, or goals from a high-level problem-space perspective~\cite{blair_models_2009, bencomo_modelsruntime_2019}.
Traditional models are commonly used for documentation, communication, and initial code generation prior to deployment.
In contrast, \textit{models@run.time} act as active abstractions during system execution.
They enable developers or the system itself to monitor, reason about, and reconfigure the software without interrupting execution~\cite{blair_models_2009}.

The fundamental mechanism that enables this continuous execution support is computational reflection, which establishes a bidirectional causal connection between the runtime model and the underlying software system.
This link keeps the internal representation aligned with the domain.
As a result, changes in the system's state or behavior are automatically captured and reflected in the runtime model.
Conversely, adaptations performed on the model at the problem-space level are propagated down to the executing system~\cite{bencomo_modelsruntime_2019}.

This causal relationship is an important prerequisite for building dynamic, self-adaptive, and long-term operating software systems.
It ensures that adaptation engines and human operators receive an accurate, up-to-date view of the system's state.
Furthermore, it allows safe adaptations to be enacted at the high level of the problem-space model, mitigating the risks associated with direct, low-level manipulations of executing code~\cite{blair_models_2009, bencomo_modelsruntime_2019}.

Despite its advantages, maintaining a continuous causal connection introduces substantial research challenges, particularly when applied to fine-grained, code-level artifacts such as abstract syntax trees~\cite{schone_incremental_2016}.
First, traditional compiler artifacts are usually discarded after compilation.
They must therefore be kept ``alive'' and enriched with runtime-specific knowledge.
Second, novel synchronization techniques are required to handle temporary synchronization lags and race conditions.
The latter occurs when multiple runtime models attempt to control the same underlying system concurrently.
Additionally, transaction-safe connections are needed to allow safe rollbacks upon failure~\cite{bencomo_modelsruntime_2019}.